% Gesamttest «Einheitskreis» — Quelle des PDFs (python3 scripts/build-lp-pdf.py einheitskreis).
% Musterlösung und Raster: bewertungspaket.tex. Alle Werte mit python3 nachgerechnet
% (scripts/lp/einheitskreis/zahlen.py, 07.10.2026; nach der Prüfung 08.10.2026 G2–G6 neu). Kein \lt/\gt — pdfLaTeX kennt sie nicht, < und > tun es.
\documentclass[11pt]{article}
\usepackage{../lp-druck}
\renewcommand{\lpname}{Einheitskreis}
\renewcommand{\lpteilgebiet}{5.4}
\renewcommand{\lpfuss}{Gesamttest · 25 Punkte}
% Einheitskreis zum Einzeichnen: Radius #1 (cm), mit Gitter alle 0.25, Achsen mit Pfeil, Zahlen ±1.
\newcommand{\ekreis}[2][]{%
\begin{tikzpicture}[scale=#2]
\draw[lplinie,line width=0.3pt] (-1.25,-1.25) grid[step=0.25] (#1);
\draw[-latex,lpgrau] (-1.3,0) -- (1.45,0) node[right,black] {$x$};
\draw[-latex,lpgrau] (0,-1.3) -- (0,1.4) node[above,black] {$y$};
\draw[black,line width=0.8pt] (0,0) circle (1);
\node[below right,font=\scriptsize,fill=white,inner sep=1pt] at (1,0) {$1$};
\node[below left,font=\scriptsize,fill=white,inner sep=1pt] at (-1,0) {$-1$};
\node[above left,font=\scriptsize,fill=white,inner sep=1pt] at (0,1) {$1$};
\node[below left,font=\scriptsize,fill=white,inner sep=1pt] at (0,-1) {$-1$};
\end{tikzpicture}}
\begin{document}
\lpkopf{Gesamttest}{%
  \vspace{4pt}\noindent\begin{tabularx}{\linewidth}{@{}X X@{}}
  Code (kein Name): \hrulefill & Punkte: \hrulefill\ / 25
  \end{tabularx}}

\begin{lpinfo}{Anleitung}
\textbf{Zeit:} rund 30 Minuten. \textbf{Hilfsmittel:} \textbf{Teil A ohne Taschenrechner} und ohne Formelsammlung, erst für \textbf{Teil B} den Taschenrechner nehmen (Modus Grad, DEG). Bleistift, Lineal und Geodreieck sind in beiden Teilen erlaubt.
Der Lehrplan vermerkt «auch ohne Hilfsmittel» bei den Beziehungen am Einheitskreis (Pythagoras, Periode, Symmetrien). Dass auch die Werte der besonderen Winkel ohne Rechner verlangt sind, ist eine Auslegung dieses Leitprogramms, wie auf der Themenseite.
Schreib zu jeder Aufgabe den \textbf{Weg} auf — eine Skizze am Einheitskreis zählt als Begründung.
Danach: Lösung scannen oder fotografieren (ohne Namen und Standort) und mit dem \emph{Bewertungspaket} bewerten lassen.
\end{lpinfo}

\lpteil{Teil A · ohne Taschenrechner}{Kapitel 1 bis 5 · 16 Punkte}

\begin{lpaufgabe}{G1}{Am Einheitskreis zeichnen}{3 P}
Im Bild ist der Einheitskreis mit der Tangente $x = 1$ gezeichnet.
\begin{enumerate}[label=(\alph*),leftmargin=*,itemsep=1pt]
\item Zeichne den Punkt $P$ zum Winkel $\varphi = 220^\circ$ ein und markiere $\sin\varphi$ und $\cos\varphi$ als Strecken.
\item Zeichne den Punkt $S$ auf der Tangente ein, dessen $y$-Koordinate $\tan 220^\circ$ ist.
\item Gib die Vorzeichen von $\sin 220^\circ$, $\cos 220^\circ$ und $\tan 220^\circ$ an.
\end{enumerate}
\par\smallskip\centering
\begin{tikzpicture}[scale=2.6]
\draw[lplinie,line width=0.3pt] (-1.25,-1.25) grid[step=0.25] (1.5,1.25);
\draw[-latex,lpgrau] (-1.3,0) -- (1.6,0) node[right,black] {$x$};
\draw[-latex,lpgrau] (0,-1.3) -- (0,1.38) node[above,black] {$y$};
\draw[black,line width=0.8pt] (0,0) circle (1);
\draw[black,line width=0.6pt] (1,-1.25) -- (1,1.25) node[above,font=\small] {$x = 1$};
\node[below right,font=\scriptsize,fill=white,inner sep=1pt] at (1,0) {$1$};
\node[below left,font=\scriptsize,fill=white,inner sep=1pt] at (-1,0) {$-1$};
\node[above left,font=\scriptsize,fill=white,inner sep=1pt] at (0,1) {$1$};
\node[below left,font=\scriptsize,fill=white,inner sep=1pt] at (0,-1) {$-1$};
\end{tikzpicture}
\par\raggedright
\lpplatz{4}
\end{lpaufgabe}

\begin{lpaufgabe}{G2}{Besondere Winkel}{4 P}
Gib exakt an (mit Wurzeln oder Brüchen). Schreib jeweils den Quadranten und den Referenzwinkel dazu — oder zeichne eine Skizze am Einheitskreis.
\begin{enumerate}[label=(\alph*),leftmargin=*,itemsep=1pt]
\item $\sin(-135^\circ)$
\item $\cos 690^\circ$
\item $\tan 120^\circ$
\item $\cos\left(-\frac{2\pi}{3}\right)$ \quad (Bogenmass)
\end{enumerate}
\lpplatz{6}
\end{lpaufgabe}

\begin{lpaufgabe}{G3}{Trigonometrischer Pythagoras}{3 P}
Es gilt $\sin\varphi = \frac{8}{17}$ und $\tan\varphi < 0$. In welchem Quadranten liegt $\varphi$? Berechne $\cos\varphi$ und $\tan\varphi$ exakt.
\lpplatz{5}
\end{lpaufgabe}

\begin{lpaufgabe}{G4}{Symmetrien}{4 P}
Es gilt $\sin 32^\circ \approx 0.530$ und $\cos 32^\circ \approx 0.848$. Gib in (a) bis (c) den Wert an und nenne jeweils die Symmetrie, die du benutzt (Spiegelung oder Formel).
\begin{enumerate}[label=(\alph*),leftmargin=*,itemsep=1pt]
\item $\sin 328^\circ$
\item $\cos 148^\circ$
\item $\cos 58^\circ$
\item Tim schreibt: $\cos 212^\circ = \cos(180^\circ + 32^\circ) = \cos 32^\circ \approx 0.848$. Was ist falsch? Gib den richtigen Wert an.
\end{enumerate}
\lpplatz{5}
\end{lpaufgabe}

\begin{lpaufgabe}{G5}{Am Einheitskreis begründen}{2 P}
\begin{enumerate}[label=(\alph*),leftmargin=*,itemsep=1pt]
\item Warum haben die Punkte zu $\varphi$ und zu $\varphi + 180^\circ$ denselben Tangens? Begründe am Einheitskreis.
\item Für welche Winkel zwischen $0^\circ$ und $360^\circ$ gibt es keinen Tangens? Begründe am Einheitskreis.
\end{enumerate}
\par\smallskip\centering\ekreis[1.25,1.25]{1.6}\par\raggedright
\lpplatz{4}
\end{lpaufgabe}

\lpteil{Teil B · mit Taschenrechner}{Kapitel 1, 4 und 5 · 9 Punkte}

\begin{lpaufgabe}{G6}{Das Schiff um die Boje}{3 P}
Ein Schiff fährt auf einem Kreis mit Radius $1$~Seemeile um eine Boje. Es startet genau östlich der Boje und fährt gegen den Uhrzeigersinn. Leg die Boje in den Ursprung: Osten zeigt in Richtung der positiven $x$-Achse, Norden in Richtung der positiven $y$-Achse. Nach einer Drehung um $235^\circ$:
\begin{enumerate}[label=(\alph*),leftmargin=*,itemsep=1pt]
\item Wie weit liegt es östlich und wie weit nördlich der Boje? (Seemeilen, drei Dezimalen; negativ heisst westlich bzw. südlich)
\item Liegt es nordöstlich, nordwestlich, südwestlich oder südöstlich der Boje?
\end{enumerate}
\lpplatz{6}
\end{lpaufgabe}

\begin{lpaufgabe}{G7}{Was der Rechner liefert}{3 P}
Lena tippt $\tan^{-1}(-2)$ und erhält $-63.4^\circ$.
\begin{enumerate}[label=(\alph*),leftmargin=*,itemsep=1pt]
\item Zeichne den Punkt $P$ zu diesem Winkel in den Einheitskreis.
\item Ein zweiter Kreispunkt hat denselben Tangenswert. Wo liegt er, und warum liefert ihn der Rechner nicht?
\item Gib für den Punkt $P$ aus (a) einen Winkel zwischen $0^\circ$ und $360^\circ$ an.
\end{enumerate}
\par\smallskip\centering\ekreis[1.25,1.25]{2}\par\raggedright
\lpplatz{3}
\end{lpaufgabe}

\begin{lpaufgabe}{G8}{Nur ein Winkel?}{3 P}
Nina rechnet $\cos^{-1}(0.7) \approx 45.6^\circ$ und schliesst: «$45.6^\circ$ ist der einzige Winkel mit $\cos\varphi = 0.7$.»
\begin{enumerate}[label=(\alph*),leftmargin=*,itemsep=1pt]
\item Stimmt das? Begründe am Einheitskreis mit einer Skizze.
\item Gib einen negativen Winkel und einen Winkel über $360^\circ$ an, die ebenfalls $\cos\varphi = 0.7$ haben.
\end{enumerate}
\par\smallskip\centering\ekreis[1.25,1.25]{2}\par\raggedright
\lpplatz{3}
\end{lpaufgabe}
\end{document}
